%%%%%%%%%%%%%%%%%%%%%%%%%%%%%%%%%%%%%%%%%%%%%%%%%%%%%%%%%%%
% Print Utility
%%%%%%%%%%%%%%%%%%%%%%%%%%%%%%%%%%%%%%%%%%%%%%%%%%%%%%%%%%%
%%%%%%%%%%%%%%%
% Conversions %
%%%%%%%%%%%%%%%
% Convert a fraction with root to a decimal number
% TODO make #3 optional
\newcommand\fractionToNum[3]{%
    \FProot\tmpA{#3}{2}
    \FPmul\tmpA\tmpA{#2}
    \FPdiv\tmpB{#1}\tmpA
    \FPset\fractionNumRes{\tmpB}
    \ensuremath{
        \tmpB
    }
}

% Convert a decimal number to a fraction with root
\newcommand\solutionC{%
    \FPset\finalmul{1}
    \FPset\tmptmptmptmp\tmptmptmp
    %
    \edef\innersequence{\xintSeq[+1]{+1}{+10}}%
    \xintFor* ##1 in \innersequence \do {%
        \FPtrunc\floored\tmptmptmptmp{0}
        \FPsub\remainder\tmptmptmptmp\floored
        \FPifgt{\remainder}{0.00001}
            \FPdiv\mulmul{1}{\remainder}
            \FPmul\finalmul\finalmul\mulmul
        \else
            \xintBreakFor
        \fi
        \FPset\tmptmptmptmp\mulmul
    }
    \FPround\finalmul\finalmul{8}
    \FPclip\finalmul\finalmul
    \FPset\finalmultmp\finalmul
    % The result will be a multiplicator that is able to round
    % the number to a whole number, but the result might not be a
    % nice number to root e.g. 8/sqrt(6), 0.09375 will result in 32
    % 0.09375 * 32 = 3, but sqrt(32) = 5,65.. but given this we know
    % that the searched for number must be a multiple of 32, in this case 64
    % Note: rounding to 6 & 36 does not work since 0.09375 * 36 = 3.375
    \xintFor* ##1 in \innersequence \do {%
        \FProot\finalmulsqrt\finalmultmp{2}
        \FPround\finalmulsqrt\finalmulsqrt{4}
        \FPclip\finalmulsqrt\finalmulsqrt
        \FPifint{\finalmulsqrt}
            \xintBreakFor
        \else
            \FPadd\finalmultmp\finalmultmp\finalmul
        \fi
    }
    \FPset\finalmul\finalmultmp
    \FProot\finalmulsqrt\finalmul{2}
    \FPround\finalmulsqrt\finalmulsqrt{4}
    \FPclip\finalmulsqrt\finalmulsqrt
    \FPmul\divisor\divisor\finalmul
    \FPmul\dividend\dividend\finalmulsqrt
}

\newcommand\numToFraction[1]{%
    \FPset\dividend{1}
    \FPset\divisor{#1}
    \FPset\minusSign{1}
    \FPiflt{\divisor}{0}
        \FPset\minusSign{-1}
    \else
    \fi
    \ensuremath{
        % reverse fraction
        \FPdiv\divisor{1}{#1}
        % Reverse root
        % Use FPmul instead of FPpow to allow for neg values
        \FPmul\divisor\divisor\divisor
        \FPset\tmptmptmp\divisor
        \FPround\divisor\divisor{10}
        \FPclip\divisor\divisor
        \FPifint{\divisor}
        \else
            \solutionC
        \fi
        %
        \FPround\divisor\divisor{4}
        \FPclip\divisor\divisor
        \FPmul\dividend\dividend\minusSign
        \FPround\dividend\dividend{4}
        \FPclip\dividend\dividend
        % TODO this should not be necessary
        \FPset\outputVar{#1}
        % Create fraction from dividend and divisor
        \FPset\outputVar{\frac{\dividend}{\sqrt{\divisor}}}
        % If the sqrt can be resolved
        \FPset\tmpDivisor{\divisor}
        \FProot\tmpDivisor\tmpDivisor{2}
        \FPround\tmpDivisor\tmpDivisor{8}
        \FPclip\tmpDivisor\tmpDivisor
        \FPifint{\tmpDivisor}
            \FPset\outputVar{\frac{\dividend}{\tmpDivisor}}
        \else
        \fi
        %
        \FPset\inputVar{#1}
        \FPmul\inputVar\inputVar\inputVar
        \FPmul\inputVar\inputVar\minusSign
        \FPround\inputVar\inputVar{8}
        \FPclip\inputVar\inputVar
        \FPifint{\inputVar}
            \FPset\outputVar{\,\sqrt{\inputVar}}
        \else
        \fi
        % If #1 was an int all along
        \FPset\inputVarA{#1}
        \FPround\inputVarA\inputVarA{10}
        \FPclip\inputVarA\inputVarA
        \FPifint{\inputVarA}
            \FPset\outputVar{\inputVarA}
        \else
        \fi
        \outputVar
    }
}

\newcommand\numToFractionA[1]{%
    \FPset\inputVarA{#1}
    \FPround\inputVarA\inputVarA{10}
    \FPclip\inputVarA\inputVarA
    % Prevent division by 0
    \FPifint{\inputVarA}
        \FPset\outputVar{\inputVarA}
        \outputVar
    \else
        \numToFraction{#1}
    \fi
}

%%%%%%%%%%%%%%%%%%%
% Scalar Creation %
%%%%%%%%%%%%%%%%%%%
% Create a scalar
% #1 out, var name to save the result (scalar) to
% #2 a, scalar value
\newcommand\newScalar[2]{%
    \expandafter\edef\csname #1\endcsname{#2}%
}

% Create a scalar as a multiple of pi
% #1 out, var name to save the result (scalar) to
% #2 a, scalar value to multiply with pi
\newcommand\newScalarPi[2]{
    \FPmul\divNumRes{\FPpi}{#2}
    \expandafter\edef\csname #1\endcsname{\divNumRes}%
}

% Create a new scalar variable with a random value
\newcommand\randomScalar[3]{
    \FPrandom\tmp
    \lerpScalar[#1]{#2}{#3}{\tmp}
}

% TODO is this useful?
% Clone the given scalar variable to a new variable
% \newcommand\cloneScalar[2][scalarCloneResult]{
% }

% Copy value from one scalar variable to another
% #1 out, var name to save the result (scalar) to
% #2 a, scalar value
\newcommand\copyScalar[2][scalarCopyResult]{
    \setScalar{#1}{#2}
}

% TODO is this useful?
% \newcommand\createScalar[1]{
% }

% Set a scalar to a value
% #1 out, var name to save the result (scalar) to
% #2 a, scalar value
\newcommand\setScalar[2]{
    \newScalar{#1}{#2}
}

% TODO decide whether to use the low level assignments or 
% use the functions like setScalar to set the variable
% Set or create a scalar with value 0
% #1 out, var name to save the result (scalar) to
\newcommand\zeroScalar[1]{
    \expandafter\edef\csname #1\endcsname{0}%
}

%%%%%%%%%%%%%%%%%%%
% Scalar Printing %
%%%%%%%%%%%%%%%%%%%
% Print the value of this scalar variable
% #1 whether or not the value should be printed as a fraction (val=1)
% #2 the value
% TODO ensure math to remove additional whitespace
\newcommand\printScalar[2][0]{%
    \FPifeq{#1}{0}%
        \roundScalar{#2}%
        \scalarRoundResult%
    \else%
        \numToFractionA{#2}%
    \fi%
}

% Print the value of this scalar variable as a multiple of Pi
% #1 whether or not the value should be printed as a fraction (val=1)
% #2 the value
\newcommand\printScalarAsPi[2][0]{%
    \FPdiv\divNumRes{#2}{\FPpi}
    \FPifeq{#1}{0}
        \roundScalar{\divNumRes}
        \scalarRoundResult
    \else
        \numToFractionA{\divNumRes}
    \fi
    \pi
}

% Print the value of this scalar variable in degree
% #1 whether or not the value should be printed as a fraction (val=1)
% #2 the value
\newcommand\printPiScalarInDegree[2][0]{%
    \FPmul\tmp{#2}{180}
    \FPdiv\divNumRes{\tmp}{\FPpi}
    \FPifeq{#1}{0}
        \roundScalar{\divNumRes}
        \scalarRoundResult
    \else
        \numToFractionA{\divNumRes}
    \fi
    ^\circ
}

% TODO rename to roundScalar / printScalarRounded
% #1 number of digits after comma, default is 4
% #2 the value
\newcommand\printScalarRounded[2][4]{%
    \FPround\tmp{#2}{#1}
    \FPclip\tmp\tmp
    \tmp
}

%%%%%%%%%%%%%%%%%%%%%%%
% Scalar Modification %
%%%%%%%%%%%%%%%%%%%%%%%
% TODO add default param scalarRoundResult does not work since
% there is already a default param
% #1 number of digits after comma, default is 4
% #2 the value
\newcommand\roundScalar[2][4]{%
    \FPround\tmp{#2}{#1}%
    \FPclip\tmp\tmp%
    \newScalar{scalarRoundResult}{\tmp}%
}

% Ceil the value of this variable
% #1 out, var name to save the result (scalar) to
% #2 the value
\newcommand\ceilScalar[2][scalarCeilResult]{
    \FPifint{#2}
    \else
        \FPadd\tmp{#2}{1}
    \fi
    \FPtrunc\tmp\tmp{0}
    \newScalar{#1}{\tmp}
}

% Floor the value of this variable
% #1 out, var name to save the result (scalar) to
% #2 the value
\newcommand\floorScalar[2][scalarFloorResult]{
    \FPtrunc\tmp{#2}{0}
    \newScalar{#1}{\tmp}
}

% TODO naming
% \newcommand\toDegree
\newcommand\radianToDegree[2][radianToDegreeResult]{
    \FPmul\tmp{#2}{180}
    \FPdiv\tmp{\tmp}{\FPpi}
    \newScalar{#1}{\tmp}
}

% TODO naming
\newcommand\toRadian[2][degreeToRadianResult]{
    \FPdiv\tmp{#2}{180}
    \FPmul\tmp{\tmp}{\FPpi}
    \newScalar{#1}{\tmp}
}

%%%%%%%%%%%%%%%
% Scalar Math %
%%%%%%%%%%%%%%%
% Add the values of two scalars
% #1 out, var name to save the result (scalar) to
% #2 a, first scalar value
% #3 b, second scalar value
\newcommand\addScalar[3][scalarAddResult]{% {Basename}{input}
    \FPadd\tmp{#2}{#3}
    \newScalar{#1}{\tmp}
}

\newcommand\subScalar[3][scalarSubResult]{% {Basename}{input}
    \FPsub\tmp{#2}{#3}
    \newScalar{#1}{\tmp}
}

% Multiply the values of two scalars
% #1 out, var name to save the result (scalar) to
% #2 a, first scalar value
% #3 b, second scalar value
\newcommand\mulScalar[3][scalarMulResult]{
    \FPmul\tmp{#2}{#3}
    \newScalar{#1}{\tmp}
}

% Divide the values of two scalars
% #1 out, var name to save the result (scalar) to
% #2 a, first scalar value
% #3 b, second scalar value
\newcommand\divScalar[3][scalarDivResult]{%
    \FPifeq{#3}{0}
        \PackageError{glmatrix package}{In divScalar function: division by zero}{parameter 3 is zero therefore you are trying to divide by 0}
    \else
    \fi
    \FPdiv\tmp{#2}{#3}
    \newScalar{#1}{\tmp}
}

\newcommand\powScalar[3][scalarPowResult]{% {Basename}{input}
    \FPpow\tmp{#2}{#3}
    \newScalar{#1}{\tmp}
}

% TODO negateScalar vs negScalar
% Negate the value of the given scalar
% #1 out, var name to save the result (scalar) to
% #2 a, scalar value
\newcommand\negateScalar[2][scalarNegateResult]{
    \FPneg\tmp{#2}
    \newScalar{#1}{\tmp}
}

% Lineraly interpolate between values of two scalars
% #1 out, var name to save the result (scalar) to
% #2 a, first scalar value
% #3 b, second scalar value
% #4 t, interpolation amount range [0-1]
\newcommand\lerpScalar[4][lerpScalarResult]{
    \FPsub\tmpA{1}{#4}
    \FPmul\tmpA{#2}\tmpA
    \FPmul\tmpB{#3}{#4}
    \FPadd\tmpA\tmpA\tmpB
    \newScalar{#1}{\tmpA}
}

\newcommand\arccosScalar[2][scalarArccosResult]{
    \FParccos\tmp{#2}
    \newScalar{#1}{\tmp}
}

% Adds two scalars after multiplying the second operand by a value
% #1 = #2 + #3 * #4
% #1 out, var name to save the result (scalar) to
% #2 a, first scalar value
% #3 b, second scalar value
% #4 c, muliplier
\newcommand\mulAndAddScalar[4][scalarMulAndAddResult]{
    \FPmul\tmp{#3}{#4}
    \FPadd\tmp{#2}{\tmp}
    \newScalar{#1}{\tmp}
}

%%%%%%%%%%%%%%%%%%%%%
% Scalar Comparison %
%%%%%%%%%%%%%%%%%%%%%
% Returns whether or not the scalars have the same values. 
% \newcommand\equalsScalar[3][scalarEqualsResult]{
% }

% Return the minimum of two scalar values
% #1 out, var name to save the result (scalar) to
% #2 a, first scalar value
% #3 b, second scalar value
\newcommand\minScalar[3][minScalarResult]{
    % TODO use \FPmin#1#2#3
    \FPiflt{#2}{#3}
        \expandafter\edef\csname #1\endcsname{#2}%
    \else
        \expandafter\edef\csname #1\endcsname{#3}%
    \fi
}

% Return the maximum of two scalar values
% #1 out, var name to save the result (scalar) to
% #2 a, first scalar value
% #3 b, second scalar value
\newcommand\maxScalar[3][maxScalarResult]{
    % TODO use \FPmax#1#2#3
    \FPiflt{#2}{#3}
        \expandafter\edef\csname #1\endcsname{#3}%
    \else
        \expandafter\edef\csname #1\endcsname{#2}%
    \fi
}

%%%%%%%%%%%%%
% Functions %
%%%%%%%%%%%%%
% Apply the given function to a list of variables.
% TODO try to use a string as #2 instead of the integer, 
% this would allow for dynamic function calls, not requireing the if else
% utelise \csname #2\endcsname for that
% #1 list of commands (var names) that should be altered
% #2 the chosen function
% #3 params for function
\newcommand\forEachScalar[3]{
    \FPiflt{#2}{10}
    \else
        \PackageError{glmatrix package}{In forEach scalar}{selection function that is not implemented (chose between 1 to 9 for parameter 2)}
    \fi
    %
    \def\varList{\xintCSVtoList{#1}}
    \def\paramList{\xintCSVtoList{#3}}
    \xintFor* ##1 in \varList \do {%
        \FPifeq{#2}{1}
            \addScalar[##1]{\csname ##1\endcsname}{\xintNthElt{1}{\paramList}}
        \else
        \fi
        \FPifeq{#2}{2}
            \subScalar[##1]{\csname ##1\endcsname}{\xintNthElt{1}{\paramList}}
        \else
        \fi
        \FPifeq{#2}{3}
            \mulScalar[##1]{\csname ##1\endcsname}{\xintNthElt{1}{\paramList}}
        \else
        \fi
        \FPifeq{#2}{4}
            \divScalar[##1]{\csname ##1\endcsname}{\xintNthElt{1}{\paramList}}
        \else
        \fi
        \FPifeq{#2}{5}
            \powScalar[##1]{\csname ##1\endcsname}{\xintNthElt{1}{\paramList}}
        \else
        \fi
        \FPifeq{#2}{6}
            \negateScalar[##1]{\csname ##1\endcsname}
        \else
        \fi
        \FPifeq{#2}{7}
            \lerpScalar[##1]{\csname ##1\endcsname}{\xintNthElt{1}{\paramList}}{\xintNthElt{2}{\paramList}}
        \else
        \fi
        \FPifeq{#2}{8}
            \mulAndAddScalar[##1]{\csname ##1\endcsname}{\xintNthElt{1}{\paramList}}{\xintNthElt{2}{\paramList}}
        \else
        \fi
        \FPifeq{#2}{9}
            \arccosScalar[##1]{\csname ##1\endcsname}
        \else
        \fi
    }
}
